% WORKING REFERENCE -- NOT PART OF THE SUBMITTED DOCUMENT.
%
% Do NOT \input this into the thesis. It is an extract from ONE paper (shen2026geometry) and is
% single-source by construction, so it can never be a thesis section: the writing skill cuts a
% section whose top citation carries more than a third of its references. Its job is to hold what
% that review establishes and what it concedes, sorted against the three things this thesis's aim
% says are needed, so the material can be spent across sections that each have their own evidence.
%
% READING DISCIPLINE. Three tiers, kept distinct below and never merged:
%   (a) Shen's own text -- read in full, docs_thesis/papers/shen2026geometry.md
%   (b) Shen's Table 7 entries -- SECOND-HAND. The primary papers have not been read. Anything
%       from the table is attributed to the table, never cited as if read.
%   (c) kashyap2022tortuosity and li2011tortuosity -- found via Table 7, then their own abstracts
%       and metadata verified directly against Europe PMC. Abstract-level, not full-text.

\section*{Shen et al.\ 2026: the field's own gaps, read against this thesis}

A narrative review of coronary hemodynamics and the anatomy driving it, by the groups behind most
of the mechanism literature, with 311 references and seven tables \cite{shen2026geometry}.
No cohort, no pooled statistic, no search protocol.
It is worth this much attention for one reason: it is the field describing its own unfinished
business, so where the thesis's aim lands on a gap the review names, the direction is the field's
and not merely ours.

\subsection*{1. The three things the aim says are needed}

\begin{description}

\item[A construction that turns a segmented tree into a measurable object --- named as a live gap.]
Shen reports that automated segmentation now allows whole coronary trees to be reconstructed at
speed, and then lists among the three things future studies must do: \emph{developing novel
algorithms to automate and accelerate coronary artery segmentation and flow analysis}.
The sharper concession is about scope rather than speed: most studies covered \emph{``the main
bifurcation regions or non-bifurcating vessels, and most curved segments and some regions commonly
susceptible to diseases were neglected''}, and they argue this matters because vulnerable plaque
sits mostly outside the left main --- spontaneous coronary artery dissection occurs in the middle to
distal segments --- so \emph{``investigating the whole tree is essential''}.
\textit{Consequence:} building the tree as a graph and measuring across all of it, rather than at
the left main, is the review's own stated requirement. It also cuts at us: the best-evidenced
feature in the state of the art is the left main bifurcation angle, and the thesis should say that
the left main is where the literature looked, not where the disease mostly is.

\item[Features a machine computes identically every time --- the review concedes the field has not
got this, and points at the fix.]
Shen sets out the definitional mess directly. Curvature is the magnitude of the derivative of the
unit tangent with respect to arc length, or intuitively $1/R$ for the osculating circle; the
tortuosity index compares path length $L$ to endpoint distance $D$ and \emph{``cannot capture the
spatial information''}. Clinical practice reduces both to a count of bends, or to a 2D projection
labeled C- or S-shaped. Computational studies have used the tortuosity index, average absolute
curvature, or a single bend, \emph{``neglecting the spatiality and continuity of this
characteristic''}. Among the named blockers on the whole enterprise is \emph{``the absence of
established thresholds for haemodynamic descriptors''}.
\textit{Consequence:} the thesis does not need to defend inventing a definition --- it needs to
defend choosing one, and the field has conceded that inconsistent measurement plausibly explains its
own contradictory findings.

\item[A cohort large enough and healthy enough --- named as the second of the review's two headline
gaps.]
\emph{``Investigations about the relationship between coronary anatomy and haemodynamics with a
large population are limited. While patient-specific geometries have increasingly been used, the
limited sample size and case selection approach may undermine the generalisability of the analyses''}.
They ask for power analysis to justify sample size, and expect future work on \emph{``more extensive
and diverse populations''}. For combining many shape factors at once they recommend statistical
shape analysis, noting it has been tested on coronary bifurcation shape
\cite{medranogracia2017bifurcation} and that \emph{``such analyses have valuable potential and are
worth further exploration''}.
\textit{Consequence:} population scale and case selection are the field's own stated limitation, and
the recommended remedy is a shape model over a population --- which is the thesis's objective 8 in
the review's vocabulary. Note the register: their ``large population'' means scaling up CFD, not
treating features as a distribution.

\end{description}

\subsection*{2. What Table 7 adds, and one thing it breaks}

Table 7 lists the reviewed in vivo and computational studies of coronary anatomy and flow across
three pages. All of the following is read from that table, not from the primary papers.

\begin{description}

\item[The scale ceiling, quantified.]
The largest whole-tree geometry-and-flow study in the table is 39 patient-specific left coronary
trees. The largest single-bifurcation CFD studies are 127 patient-specific left main bifurcations.
Only the in vivo tortuosity-scoring studies reach the order of a thousand patients, and they define
tortuosity by counting bends on an angiogram.
\textit{Consequence:} this is the scale gap in one sentence, sourced from the field's own summary
table --- whole-tree geometry has never been measured above $n = 39$, and the only measurements at
$n \approx 1000$ use the crudest available definition.

\item[Bifurcation angle: the simulations do not agree with the clinic. \textbf{This contradicts a
sentence the thesis currently asserts.}]
The clinical direction in Table 7 is consistent --- wider angle with more disease --- and includes
Temov et al.\ at 196 patients \cite{temov2016bifurcation}, Cui et al.\ at 106
\cite{cui2017bifurcation}, and Sun and Cao at 30, who report $94^\circ \pm 19.7^\circ$ in diseased
against $75.5^\circ \pm 19.8^\circ$ in non-diseased left coronary arteries ($p = 0.02$).
The computational direction is not consistent. Beier et al.\ (2016) found the angle showed
\emph{no significant effect} on adverse hemodynamics; Chiastra et al.\ (2017) found a
\emph{minor impact} in both stenosed and unstenosed models; Shen et al.\ (2021) found that a
\emph{larger angle reduced} the area exposed to low shear, which is the opposite of the expected
direction. Others in the same table point the expected way.
\textit{Consequence:} \texttt{state\_of\_the\_art/05\_predicts\_disease.tex} currently closes its
bifurcation-angle argument with ``an anatomical gradient, an outcome association and a hemodynamic
mechanism pointing the same way''. On the strength of one CFD study
\cite{wang2024lmphenotypes}, that sentence overstates a literature the anchor review shows to be
split, and it needs rewriting.

\item[Tortuosity: the sign depends on the level of analysis.]
At artery and patient level the association is negative or null. Li et al.\ scored 1010 consecutive
angiography patients and found tortuosity \emph{negatively} correlated with coronary artery disease,
odds ratio 0.755, 95\% confidence interval 0.574 to 0.994, $p = 0.045$, with major adverse events
over two to four years no different between diseased patients with and without it
\cite{li2011tortuosity}. Chiha et al.\ at 870 patients found women with severe tortuosity more
likely to have normal arteries. At segment level it inverts: Tuncay et al.\ at 73 patients report
tortuosity 30\% higher for segments carrying plaque, \emph{with no significant correlation at the
artery level}. Several CFD entries propose the mechanism --- tortuosity induces helical flow, which
is offered as atheroprotective.
\textit{Consequence:} tortuosity has no single sign, and the thesis must state the level of analysis
before reporting any tortuosity result. The segment-level and patient-level findings are not in
conflict; they are different measurements wearing one word.

\item[Sex is a systematic modifier, not a covariate to adjust away.]
Across four independent entries: males were 2.07-fold more likely to exceed an $80^\circ$
bifurcation angle \cite{temov2016bifurcation}; tortuosity was markedly more common in women
(odds ratio 2.603, 95\% confidence interval 1.897 to 3.607, $p < 0.001$) \cite{li2011tortuosity};
women with severe tortuosity were more likely to have normal arteries; and a 127-case CFD study
found low shear on the \emph{inner} wall of curvature in females and the \emph{outer} wall in males,
with more adverse oscillatory shear in males.
\textit{Consequence:} the last of these is a difference in \emph{where} the adverse shear sits, not
how much of it there is. Every distribution this thesis reports is stratified by sex, and the
justification is this cluster rather than convention.

\end{description}

\subsection*{3. The measure to use, and why}

Found through Table 7 and then verified directly, so this is the one entry here resting on its own
source rather than on the review. Kashyap et al.\ took CTCA from \textbf{127 patients without coronary
artery disease}, applied every previously used tortuosity metric to the left main bifurcation and
its two branches, and correlated each against the percentage of vessel area below 0.4 Pa
time-averaged wall shear stress \cite{kashyap2022tortuosity}.
Average absolute curvature had the highest coefficient of determination across all branches
($p < 0.001$), ahead of average squared-derivative curvature ($p = 0.001$) and root-mean-squared
curvature ($p = 0.002$).
\textbf{The tortuosity index --- the measure most widely used in the literature --- showed no
significant correlation with low shear at all ($p = 0.86$).}
\textit{Consequence:} objective 7 computes average absolute curvature, and says why. It also
qualifies the earlier conclusion in \texttt{find\_research.md} \S1.2: a continuous measure beats the
clinical bend count against a clinical endpoint, but among continuous measures the tortuosity index
is the wrong one against a shear endpoint.

\subsection*{4. What this review does not license}

It is a narrative review with no search protocol, no inclusion criteria and no screening counts, so
\emph{``the field has not done X''} cannot be sourced to it --- only \emph{``the field says it has
not done X''}.
It carries no cohort and no pooled estimate, so no number in it can be quoted with a population
behind it except by going to the primary paper.
The words \emph{outlier}, \emph{extreme value} and \emph{anomaly} appear nowhere in it: its
population framing is about running more and better simulations on larger cohorts, not about
treating features as a distribution and flagging its tails. Objective 9 gets no endorsement here,
and claiming one would be caught by anyone who opens the paper.
